\section{An explicit Eisenstein eightfold bound after Schoen and Severn}
\label{sec:eisenstein-eightfolds}
This section bounds the degrees in the cyclic-cover transfer used for the
split Eisenstein eightfold construction. The geometric input is the
$G_{34}$ family of Severn's unrefereed preprint
\cite[Proposition~8.5, \S9 and Appendix~A]{Severn}. The calculation below
is a quantitative consequence of that construction. It does not use
Perry's theorem or either CM Hodge-conjecture manuscript.

Put
\begin{align*}
 G&=39191040,& g&=4G+1,\\
 R&=(2G+1)G!G^G,& L_0&=2gR,\\
 S_0&=G^{2g}(g!)^2,& S_B&=(L_0/2)!S_0^R,\\
 k&=L_0/2+4,&
 B_0&=(40G)^k\binom{k}{4}S_B\,8!G^8.
\end{align*}
These formulas specify integers; they do not require decimal expansion.
The constant is deliberately coarse, so that no estimate for coefficients
in an invariant-theory decomposition enters the proof.

\begin{theorem}\label{thm:eisenstein-eightfolds}
Let $K=\Q(\sqrt{-3})$, let $\mathcal O_f=\Z+f\mathcal O_K$ act
integrally on a polarized abelian variety of Weil type $(A,L_A)$, and
assume that the polarization is compatible with the action. Write
$\Delta=\chi(L_A)$. Using the $G_{34}$ geometric input just specified,
every split Weil eightfold satisfies
\[
 \delta_W(A,L_A)\le B_0 f^4\Delta.
\]
Every Weil sixfold over $K$, of arbitrary determinant, satisfies
\[
 \delta_W(A,L_A)\le \frac{B_0}{4}
          \frac{f^9\Delta^2}{b},
 \qquad b=[\mathcal O_K H_1(A,\Z):H_1(A,\Z)].
\]
The input from the preprint is a specific dominant split family of
group-algebra factors of Jacobians. The proof gives effective integral
detecting cycles and all the degree losses from that family to $A$.
\end{theorem}

\subsection{The geometric input and the factor lattice}

Here is the precise form of the geometric input. The complex reflection
group $\mathcal G=G_{34}$ has order $G$, has its faithful six-dimensional
reflection representation over $K$, and contains the scalar $\omega I$.
Ten involutory reflections $s_1,\ldots,s_{10}$ generate it. The tuple
\[
 (s_1,s_1,\ldots,s_{10},s_{10})
\]
defines a Hurwitz family of connected $\mathcal G$-covers
$C\to\mathbf P^1$. Its reflection multiplicity in $H^1(C,K)$ has
$K$-dimension eight and Hodge numbers $(4,4)$. The period map of the
resulting abelian eightfold factor has an open image in the split Weil
domain. These are the statements proved in \cite[\S9]{Severn}; the
matrices and the seventeen boundary twists are in its Appendix~A.
Riemann--Hurwitz gives $2g(C)-2=8G$, explaining the integer $g$ above.

We recall why the period certificate has the stated consequence. On the
rank-eight parabolic cohomology, the commuting logarithms $N_j$ have
$N_jN_l=0$, their images span a four-dimensional isotropic $K$-space
$I$, and they annihilate $I$. Their associated Hermitian forms on the
quotient are semidefinite of a common sign, with definite sum. They span
a rational space of dimension sixteen. Geometric monodromy logarithms
are common-sign positive multiples of these logarithms after a finite
base change. A positive sum therefore has image and kernel $I$. The
limiting mixed Hodge structure has weight-zero part $I$ of type $(0,0)$,
so its limiting $F^1$ is complementary to $I$. The rational Lie algebra
of unitary maps from the quotient to $I$ has dimension sixteen; its
complexification maps isomorphically to the tangent space at that
limiting filtration. Thus sixteen of the $N_j$ give independent tangent
vectors. Holomorphic untwisting of the period map shows that the same
minor is nonzero at nearby smooth fibres. The isotropic space $I$ also
proves that the Hermitian space is split. This is the implication in
\cite[Lemmas~8.2--8.4 and Proposition~8.5]{Severn}; its finite matrix
input, including the nonzero minor $1728$, is distinct from the degree
argument that follows.

Let $\chi$ be the character of the reflection representation, and let
$s$ be a reflection. In the rational group algebra put
\[
 \varepsilon=\frac6G\sum_{a\in\mathcal G}
       \operatorname{Tr}_{K/\Q}(\chi(a^{-1}))a,
 \qquad e=\varepsilon(1-s)/2.
\]
The first operator is the central projector for the conjugate pair of
reflection representations. The second has rank one in its matrix
algebra $M_6(K)$. Its image on $J(C)$ is an abelian eightfold, denoted
$A_C$, with the restricted polarization $\theta_C$. Character values
are algebraic integers, so their rational traces are integers and
$Ge\in\Z[\mathcal G]$. The Rosati involution sends $a$ to $a^{-1}$;
both commuting factors of $e$ are self-adjoint. The scalar $\omega I$
preserves $A_C$, giving integral $\mathcal O_K$-action.

Let $\Lambda$ be the unimodular lattice of $J(C)$ and
$\Lambda_C=\Lambda\cap e\Lambda_\Q$. This sublattice is primitive.
Every integral functional on it extends to $\Lambda$, and
self-adjointness identifies its dual lattice with $e\Lambda$.
Consequently
\[
 \Lambda_C^{\#}=e\Lambda,\qquad
 G\Lambda_C^{\#}\subset\Lambda_C.
\]
Every elementary divisor of the restricted polarization divides $G$.
In particular
\begin{equation}\label{eq:eis-source-volume}
 \chi(\theta_C)\mid G^8,
 \qquad \int_{A_C}\theta_C^8\le8!G^8.
\end{equation}

\subsection{A bounded integral relation between determinant characters}

Let $U$ be the full centralizer of $\mathcal G$ in the symplectic group
of $H^1(C,\Q)$. Over $\C$, a nonreal conjugate pair of irreducible
representations gives a general linear multiplicity group, an orthogonal
irreducible representation gives a symplectic multiplicity group, and
a quaternionic one gives an orthogonal multiplicity group. Hence the
character group of $U$ has a free determinant coordinate for each
nonreal pair and a coordinate of order two for each quaternionic
multiplicity group. The full orthogonal groups must be used here.

For a subgroup $H'\le\mathcal G$ and a linear character $\chi'$ of
$H'$, let $a(\chi')$ be the determinant character of its multiplicity
space. Frobenius reciprocity identifies this with the image of
$\operatorname{Ind}_{H'}^{\mathcal G}\chi'$ under the homomorphism
which takes differences of conjugate multiplicities and quaternionic
multiplicities modulo two. Brauer induction, in the form used in
\cite[Theorem~5.1]{Severn}, gives an integral relation
\begin{equation}\label{eq:eis-supplier-relation}
 e_\sigma=\sum_i n_i a(\chi_i),
\end{equation}
where $e_\sigma$ is the character of one target Weil line.

Characters of order one or two contribute zero. Their multiplicities
in conjugate nonreal representations agree; their multiplicities in
a quaternionic representation are even, since their isotypic subspace
inherits a nondegenerate alternating form. A character with zero
cohomological multiplicity also contributes zero. All remaining terms
come from cyclic intermediate covers
\[
 C_i'=C/\ker\chi_i\longrightarrow X_i=C/H_i'
\]
of order $m_i\ge3$. Since the original inertia groups have order two,
every nonzero branch exponent in such a cyclic cover is $m_i/2$.
The sum of the exponents is zero modulo $m_i$, so their number $s_i$ is
even and they can be paired with their negatives. If $q_i=g(X_i)$,
the primitive eigenspaces in $H^1(C_i')$ have common dimension
$h_i=2q_i-2+s_i$. A nonzero retained multiplicity has $h_i>0$, and
$h_i$ is even. Unramified intermediate covers are included.

There is a uniform bound
\begin{equation}\label{eq:eis-relation-mass}
 \sum_i|n_i|\le R.
\end{equation}
To prove it, append columns $2e_j$ for the order-two coordinates and
lift the relation to an integer column lattice. There are at most $G$
rows, and the column entries have absolute value at most $G$. Restrict
to independent rows; this does not change relations between the
columns. If their rank is $r$, choose an invertible integer matrix $B$
of $r$ columns and put $d_0=|\det B|\le r!G^r$. In the subgroup of
$\Z^r/B\Z^r$ generated by all columns, a shortest expression for the
target uses at most $d_0-1$ signed columns: a shortest path has no
repeated vertex. The residual vector has coordinates at most $d_0G$.
Cramer's rule, with cofactors at most $(r-1)!G^{r-1}$, expresses it in
the columns of $B$ with each coefficient at most $r!G^r$. The total
coefficient sum is therefore at most $(r+1)r!G^r\le R$. Removing the
auxiliary parity columns leaves \eqref{eq:eis-supplier-relation} and
cannot increase that sum.

\subsection{Integral cycles from the cyclic intermediate covers}

Fix one term, and write $C'$, $X$, $m$, $h$ for its data. We use
Schoen's construction \cite[Theorem~2.0 and Lemmas~2.3--2.6,
pp.~11--17]{Schoen88}. The primitive rational part has dimension
$\varphi(m)h$, so
\[
 h\le g(C')\le g,
 \qquad m\le G.
\]
Let $G_0$ be the kernel of the sum map on $(\Z/m)^h$, extended by
the permutation group on the $h$ factors. It has order
$m^{h-1}h!$. There are finite quotient maps
\[
 q:C'^h\longrightarrow W_0=C'^h/G_0
       \longrightarrow\operatorname{Sym}^hX,
\]
the second with residual cyclic group $\Z/m$. Schoen chooses a linear
$P_0\simeq\mathbf P^{h/2}$ in a complete linear system on $X$. The
strict transform of $P_0$ is unchanged by a sequence of blowups and
avoids the branch divisors. The cyclic cover therefore splits over it.
A component descends to $Q_0\subset W_0$ birationally over $P_0$,
and its primitive-character components in cohomology are nonzero.

We specify the exceptional-divisor issue in this last assertion.
When $4\mid m$ and all exponents are $m/2$, the prime chosen in
\cite[Sublemma~2.9, p.~19]{Schoen88} does not exist. The replacement
is the faithful-character argument of \cite[Lemma~4.4, pp.~9--10]{Severn}.
At a blowup of the intersection of a paired pair of branch divisors,
their total exponent is zero modulo $m$. The exceptional divisor is a
$\mathbf P^1$-bundle and is not itself a branch divisor. The two
strict transforms meet each fibre in two points. For a faithful
character, the eigensheaf is extension by zero of a rank-one local
system on their complement $\C^\times$, with nontrivial monodromy.
Its compactly supported cohomology vanishes: degree zero is zero,
degree two is dual to zero invariants, and the Euler characteristic is
zero. If a whole fibre lies in another branch divisor, its eigensheaf
is zero because its inertia is nontrivial. Proper base change thus
kills the faithful-character cohomology of the entire exceptional
divisor. The original centre has nontrivial inertia as well and has
zero faithful-character eigensheaf. The relative-cohomology sequence
then identifies the primitive parts before and after the blowup.
Schoen's nonzero normal-bundle intersection on the final component
therefore remains nonzero on $Q_0$. This proves the needed descent for
these exponents, including those omitted by the prime-choice argument.

Let $H$ denote the sum of point classes on $C'^h$. An incidence divisor
$D_x=x+\operatorname{Sym}^{h-1}X$ restricts to the hyperplane class on
$P_0$, and pulls back to $mH$ on $C'^h$. If $Z_0$ is an irreducible
component above $Q_0$ and has degree $e_0$ over $Q_0$, then
\[
 e_0\le |G_0|,
 \qquad \deg_H Z_0=e_0/m^{h/2}.
\]
The integral orbit sum
$\mathcal R_0=\sum_{a\in G_0}a_*Z_0$, with repetitions, has effective
degree at most $|G_0|^2/m^{h/2}$. Rational transfer for a finite
quotient gives $[\mathcal R_0]=e_0q^*[Q_0]$. This cohomological
identity does not require the quotient to be flat along $Q_0$.

For the residual action $\delta$ of $\Z/m$, form the integral signed
cycle
\[
 Y=\sum_{t=0}^{m-1}c_m(t)\delta(t)_*\mathcal R_0,
\]
where $c_m(t)$ is the Ramanujan sum. This operator is $m$ times the
rational primitive projector, and $|c_m(t)|\le\varphi(m)$. It keeps
the nonzero primitive components. In $G_0$-invariants a primitive
residual character forces all coordinate characters to be the same
primitive character. Consequently only
$H^1(C')^{\otimes h}$ occurs, and permutation invariance identifies
its component with the corresponding determinant line; the signs in
the permutation action are the cohomological signs in degree one.
Writing the mass of a signed cycle for the sum of the degrees of its
positive and negative parts, we obtain
\begin{equation}\label{eq:eis-cyclic-mass}
 \operatorname{mass}_H(Y)
 \le m\varphi(m)\frac{(m^{h-1}h!)^2}{m^{h/2}}
 \le m^{3h/2}(h!)^2\le S_0.
\end{equation}
The orbit sum in this estimate includes its integral multiplicities.

\subsection{Transfer by a finite collection of divisors}

Take $|n_i|$ copies of the cycle for each term of
\eqref{eq:eis-supplier-relation}. Their exterior product $Y_B$ lies on
a product $B$ of $\ell=\sum_i|n_i|h_i$ curves, with
$\ell\le gR\le L_0$. For the sum $H_B$ of point classes,
\begin{equation}\label{eq:eis-product-mass}
 \operatorname{mass}_{H_B}(Y_B)
 \le \frac{(\ell/2)!}{\prod_i((h_i/2)!)^{|n_i|}}S_0^R
 \le S_B.
\end{equation}
Every selected determinant component is nonzero. The perfect pairing
between inverse primitive eigenspaces on each curve gives its dual
determinant tensor. Choose the component according to the sign of
$n_i$. The product of the dual tensors with one target Weil eigenline
has degree $\ell+8$ on $B\times A_C$, is $U$-invariant by
\eqref{eq:eis-supplier-relation}, and sends $[Y_B]$ to a nonzero vector
of that line. Denote this tensor by $\Gamma$. Both the source cycles
and their dual tensors have middle degree on the curve products.

We need control of the divisors representing $\Gamma$. Let
$D$ be the image of $\Q[\mathcal G]$ in $\operatorname{End}H^1(C)$.
The semisimple-algebra lemma and the double-centralizer argument in
\cite[Lemma~3.5]{MilneLefschetz} give
$\operatorname{End}_U H^1(C)=D$. The invariant theorem for full
symplectic and orthogonal groups and paired general linear groups
\cite[Proposition~3.6 and Remark~3.7]{MilneLefschetz} says that their
exterior invariants on powers are generated in degree two.
The degree-two tensors correspond to self-adjoint matrices over $D$.
Off-diagonal entries and their adjoints are spanned by group elements;
the diagonal self-adjoint entries are spanned by $a+a^{-1}$.
It follows that these tensors are spanned by the divisor classes of
the globally generated line bundles
\begin{equation}\label{eq:eis-divisor-list}
 \operatorname{pr}_i^*(2\Theta),\qquad
 (\operatorname{pr}_i+a\operatorname{pr}_j)^*(2\Theta),
 \qquad a\in\mathcal G,
\end{equation}
on powers of $J(C)$. Repeated indices are allowed: the identity
$(1+a)^\dagger(1+a)=2+a+a^{-1}$ supplies the diagonal entries.

For a quotient $\pi:C\to C'$, the map
$C'\to J(C')\xrightarrow{\pi^*}J(C)$ identifies its first cohomology
with a rational direct summand, equivariantly for $U$. The idempotent
$e$ gives the corresponding summand for $A_C$. Thus $\Gamma$ lifts
to an exterior invariant on a power of $J(C)$. Its degree is $\ell+8$,
so it is a complex linear combination of monomials of
$c=\ell/2+4$ divisors from \eqref{eq:eis-divisor-list}. Since its
action on $Y_B$ is nonzero, some monomial has nonzero projection to
the target Weil line. No bound for the coefficients of the linear
combination is required.

We now bound every such monomial. If $a_0=\deg\pi$ and
$\eta:C'\to J(C)$ is the preceding map, the norm identity and
Riemann--Hurwitz give
\[
 \deg\eta^*\Theta=a_0g(C')
 \le g(C)-1+a_0\le5G.
\]
On $B\times A_C$, the restricted product polarization is therefore
at most $5G(H_B+\theta_C)$ in the nef order. The theorem of the
square bounds each line bundle in \eqref{eq:eis-divisor-list} by eight
times that product polarization, also for repeated indices. Hence
\[
 D_j|_{B\times A_C}\le40G(H_B+\theta_C)
\]
in the nef order. Successive general members of these globally
generated systems meet every component of $|Y_B|\times A_C$
properly. They give effective integral four-cycles. Their signed
pushforward to $A_C$ has the prescribed correspondence class, and
its degree mass is at most
\begin{align*}
 &(40G)^c\int_{|Y_B|\times A_C}
                    (H_B+\theta_C)^c\theta_C^4\\
 &\qquad=(40G)^c\binom c4
              \operatorname{mass}_{H_B}(Y_B)
              \int_{A_C}\theta_C^8
 \le B_0.
\end{align*}
Only the term with exponents $(\ell/2,4)$ contributes. Components
whose images have smaller dimension push to zero. Since the total
Weil projection is nonzero, one irreducible component of the remaining
effective cycles detects. This proves $\delta_W(A_C,\theta_C)\le B_0$.

\subsection{From the source to every polarization and determinant}

The source has maximal-order Eisenstein action and split Hermitian
space. Since $t_K=1$ and the class number is one,
Theorem~\ref{thm:principal-quotient} gives an isogeny to the principal
$U^4$ family pulling the target polarization back to $\theta_C$.
Pushing forward a detecting cycle costs at most its original degree.
The induced map between rational period domains is an isomorphism,
so the principal family still has an open image with detector bound
$B_0$. The closed bounded-detector locus of
Proposition~\ref{prop:detector-loci} extends this bound to its entire
connected component. Pullback through the principal quotient of an
arbitrary target costs $f^4\Delta$, proving the eightfold estimate.

For a maximal-order sixfold of arbitrary determinant, the balanced CM
surface in Theorem~\ref{thm:explicit-stabilization} has polarization
product $\Delta$. Its product with the sixfold is a split eightfold
of polarization product $\Delta^2$. The opposite-graph contraction
in that theorem has cost $1/4$ in the Eisenstein case. This gives
$B_0\Delta^2/4$. For an order of conductor $f$, saturation has degree
$b$, polarization multiplier $f$, and polarization product
$\Delta'=f^6\Delta/b$. Pulling back a three-cycle costs $b/f^3$.
Thus the bound is
\[
 \frac b{f^3}\frac{B_0}{4}
       \left(\frac{f^6\Delta}{b}\right)^2
 =\frac{B_0}{4}\frac{f^9\Delta^2}{b},
\]
as asserted. The factor $8!$ in $B_0$ makes $B_0/4$ an integer.
